\section{Conclusion}
\label{sec:conclusion}

Evaluated on floods the model has never seen, terrain helps a SAR flood-segmentation network as an input:
a U-Net trained on the 11 Sen1Floods11 events improves water and flood mapping on 79 external CEMS
activations. Its event-level part needs correctly located, detailed height above drainage, which relief, a
permanent-water map or, with few training events, a larger pretrained model does not replace.
As a post-processing mask HAND recovers part of the gain, as a physics loss little, and with 159 training
events its flood-extent gain is not detectable. Whether same-event splits understate the gain, as the
Bolivia hold-out suggests, would need more than Sen1Floods11's ten test events to settle. With few labelled
floods, terrain is a cheap gain, and it should be judged on held-out floods.

\section*{Data and Code Availability}
All data are public: Sen1Floods11~\cite{bonafilia2020sen1floods11}, GEOID-Flood~\cite{geoidflood} (CC BY
4.0), Copernicus GLO-30~\cite{glo30dem}, ASF GLO-30 HAND~\cite{asfhand}, ESA WorldCover~\cite{zanaga2022worldcover}
and JRC Global Surface Water~\cite{pekel2016gsw}. Code (MIT license), every SLURM submission and the scripts
that regenerate every table are at \url{https://github.com/jeremyky/SAR-Flood-Terrain}; run outputs and
checkpoints are available from the authors and will be archived with a DOI on publication.

\section*{Acknowledgment}
The authors acknowledge Research Computing at the University of Virginia for providing computational
resources and technical support that have contributed to the results reported in this publication.
